\section{Aqueous Reactions and Gases}

\subsection{Solutions and Electrolytes}

\subsubsection{Molarity}
\[
c=\frac{n}{V}
\]
Molarity is amount of substance per solution volume. Formal concentration \(c\) is distinguished from actual species concentrations written with brackets.

\subsubsection{Actual ion concentrations}
\[
\ce{CaCl2(aq) -> Ca^2+(aq) + 2Cl-(aq)}, \qquad
[\ce{Ca^2+}]=c,\quad [\ce{Cl-}]=2c
\]
Strong electrolytes are treated as fully dissociated in water. Bracket notation gives actual concentration of each dissolved species.

\subsubsection{Strong and weak electrolytes}
\[
\ce{HCl(aq) -> H+(aq) + Cl-(aq)}, \qquad
\ce{CH3COOH(aq) <=> H+(aq) + CH3COO-(aq)}
\]
Strong electrolytes ionize essentially completely. Weak electrolytes establish equilibria with both molecular and ionic forms present.

\subsection{Precipitation and Net Ionic Reactions}

\subsubsection{Precipitation reaction}
\[
\ce{Pb(NO3)2(aq) + 2KI(aq) -> PbI2(s) + 2KNO3(aq)}
\]
Precipitation occurs when ions form a sparingly soluble solid. Spectator ions can be removed to obtain the net ionic reaction.

\subsubsection{Net ionic equation}
\[
\ce{Pb^2+(aq) + 2I-(aq) -> PbI2(s)}
\]
The net ionic equation shows only species that chemically change. It is the compact form used for precipitation and acid--base reactions.

\subsection{Redox in Aqueous Solution}

\subsubsection{Redox half-reactions}
\[
\ce{Mg -> Mg^2+ + 2e-}, \qquad
\ce{1/2O2 + 2e- -> O^2-}
\]
Oxidation produces electrons and reduction consumes electrons. The electron count must cancel in the summed redox reaction.

\subsubsection{Balancing redox reactions}
\[
\begin{aligned}
&\text{balance atoms except H,O},\\
&\text{balance O with }\ce{H2O},\\
&\text{balance H with }\ce{H+}\text{ or }\ce{OH-},\\
&\text{balance charge with }\ce{e-}.
\end{aligned}
\]
In acidic solution use \(\ce{H+}\); in basic solution neutralize added \(\ce{H+}\) with \(\ce{OH-}\). The final equation must conserve atoms and charge.

\subsection{Gases}

\subsubsection{Ideal gas law}
\[
pV=nRT
\]
The ideal gas law relates pressure, volume, amount, and absolute temperature. It is used for gas stoichiometry and pressure changes.

\subsubsection{Combined gas law}
\[
\frac{p_1V_1}{T_1}=\frac{p_2V_2}{T_2}
\]
For fixed \(n\), pressure-volume-temperature changes follow this ratio. Temperatures must be in kelvin.

\subsubsection{Partial pressure and mole fraction}
\[
p_i=x_i p_{\mathrm{tot}}, \qquad x_i=\frac{n_i}{n_{\mathrm{tot}}}, \qquad
p_{\mathrm{tot}}=\sum_i p_i
\]
Dalton's law expresses total pressure as the sum of component partial pressures.

\subsubsection{Gas density and molar mass}
\[
\rho=\frac{m}{V}=\frac{pM}{RT}, \qquad M=\frac{\rho RT}{p}
\]
Gas density can determine molar mass when the gas behaves ideally.

\subsubsection{Graham's law of diffusion}
\[
\frac{r_1}{r_2}=\sqrt{\frac{M_2}{M_1}}
\]
At equal temperature, lighter gases diffuse or effuse faster than heavier gases.
